\section{Introduction}

The previous chapter positioned AI-Scientist within the literature on grounded generation, claim verification, self-critique, and multi-agent orchestration \cite{lewis2020rag,thorne2018fever,shinn2023reflexion,wu2023autogen}. This chapter converts that background into a precise problem definition. The aim is to define the system in a way that is both technically useful and thesis-ready: every important concept should be measurable, every verdict should be interpretable, and every design choice should follow naturally from the research problem.

AI-Scientist is defined here as a verification-oriented research system that receives a scientific question, retrieves layered evidence, decomposes that evidence into checkable claims, assigns a verdict and confidence to each claim, and produces a traceable report \cite{petroni2020kilt,thakur2021beir,es2023ragas}. That framing makes the thesis stronger because it turns a broad idea into a well-scoped and academically defensible system contribution.

\begin{figure}[H]
\centering
\begin{tikzpicture}[
    node distance=1.0cm,
    scale=0.95,
    transform shape,
    box/.style={rectangle, rounded corners, draw=black, fill=blue!18, align=center, minimum width=4.3cm, minimum height=1.0cm},
    arrow/.style={-{Latex[length=3mm]}, thick}
]
\node[box] (q) {Research question};
\node[box, below=of q] (r) {Claim extraction\\and verification};
\node[box, below=of r] (c) {Critique,\\synthesis, and\\report};
\draw[arrow] (q) -- (r);
\draw[arrow] (r) -- (c);
\end{tikzpicture}
\caption{Problem-definition pipeline for AI-Scientist}
\label{fig:problem_definition_pipeline}
\end{figure}

\section{Task Definition}

The task addressed in this thesis is the following: given a natural-language research question \(q\), produce a structured and grounded response using a multi-layer evidence environment \cite{guu2020realm,karpukhin2020dpr,khattab2020colbert}. The evidence environment may include user-uploaded papers, a curated offline corpus, and live scholarly sources when available. The output is a structured verification report that reflects how strongly each claim is supported by evidence, together with a concise answer that the user can read immediately after the verified claims.

This task is intentionally broader than a single retrieval setting and narrower than unrestricted generation. It is broad because it must work across domains and question types. It is focused because the system is constrained to research-oriented queries and to evidence-based verdicts. That balance makes the problem realistic and academically defensible, and it is consistent with the design direction seen in retrieval-augmented and knowledge-intensive language systems \cite{lewis2020rag,petroni2020kilt,borgeaud2022retro}.

The system input and output can be summarized as follows, following the same kind of structured specification used in benchmark-style QA and verification tasks \cite{rajpurkar2016squad,joshi2017triviaqa,yang2018hotpotqa}:
\begin{table}[H]
\centering
\begin{tabularx}{\linewidth}{>{\raggedright\arraybackslash}p{0.24\linewidth} >{\raggedright\arraybackslash}X}
\toprule
Element & Description \\
\midrule
Input question & A natural-language scientific or technical question submitted by the user \\
User evidence & Optional uploaded papers for the current session \\
Offline corpus & Curated papers that are always available for grounding \\
Live sources & Optional retrieval from scholarly APIs for broader coverage \\
Output & A structured report with claim verdicts, confidence, evidence snippets, and rationale \\
\bottomrule
\end{tabularx}
\caption{Input and output specification for AI-Scientist}
\end{table}

\begin{table}[H]
\centering
\begin{tabularx}{\linewidth}{>{\raggedright\arraybackslash}p{0.23\linewidth} >{\raggedright\arraybackslash}X}
\toprule
Verdict & Interpretation \\
\midrule
Supported & The evidence clearly aligns with the claim and provides a stable basis for the response \\
Contradicted & The evidence points in the opposite direction or introduces a direct conflict \\
Insufficient evidence & The available material is promising but not yet strong enough to justify a firm claim \\
Out of scope & The input falls outside the research-focused operating range of the system \\
\bottomrule
\end{tabularx}
\caption{Verdict interpretation used in the thesis}
\end{table}

\section{Verdict Space and Claim-Level Reasoning}

A central decision in the thesis is to make the claim, not the whole answer, the unit of reasoning \cite{honovich2022true,min2023factscore}. This is a better fit for scientific work because papers rarely support every detail of an answer equally. Some parts may be strongly supported, some may be weakly supported, and some may remain unresolved. A single answer-level label would hide these differences, whereas a claim-level verdict makes them visible and useful.

AI-Scientist therefore uses a bounded verdict space:
\begin{itemize}[leftmargin=*]
    \item \textbf{Supported} when the retrieved evidence clearly aligns with the claim,
    \item \textbf{Contradicted} when the evidence supports the opposite direction or clearly conflicts with the claim,
    \item \textbf{Insufficient evidence} when the available evidence is too weak to decide confidently,
    \item \textbf{Out of scope} when the input is not a research question and should be declined early.
\end{itemize}

This verdict space has an important thesis value: it encourages clarity without reducing the utility of the system. In other words, the system remains effective when it chooses a careful verdict. That is a strong research property because it shows that the framework knows not only how to answer, but also how to calibrate the strength of its response to the underlying evidence.

\section{Evidence Model and Relevance}

The quality of a verdict depends on the quality of evidence behind it. For this reason, AI-Scientist uses a layered evidence model rather than a flat corpus. The evidence layers are ordered by practical trust and relevance:
\begin{enumerate}[leftmargin=*]
    \item user-uploaded papers supplied for the session,
    \item curated offline corpus papers,
    \item live scholarly sources retrieved during execution.
\end{enumerate}

This structure is useful because it mirrors the way a careful researcher works. User-provided material should be respected when it is relevant, the curated corpus provides a dependable grounding baseline, and live retrieval broadens the context when extra support is needed. Source priority is therefore not a side detail; it is part of the problem definition itself, and it follows the same practical layering logic used in retrieval-oriented pipelines \cite{ragsurvey2024,es2023ragas}.

To make evidence selection transparent, the thesis uses a coverage-based relevance measure for papers and a token-overlap measure for finer sentence-level support. A paper is considered admissible when it clears a minimum relevance threshold, which helps maintain a focused and reliable verification chain while keeping the retrieval stage efficient.

\begin{equation}
\mathrm{cov}(q, d) = \frac{|T(q) \cap T(d)|}{|T(q)|}
\end{equation}

Here, \(T(q)\) and \(T(d)\) are the normalized token sets of the query and document. This formulation is attractive because it rewards a document for covering the query terms without unfairly penalizing longer abstracts. In practice, it gives the pipeline a concise first indicator of whether a paper is worth deeper inspection.

\section{Confidence and Traceability}

Confidence is treated as a first-class output in this thesis \cite{kadavath2022language,lin2021truthfulqa}. That choice is important because a scientific assistant should not only say what it believes; it should also indicate how strongly the evidence supports that belief. The confidence score combines evidence overlap, corroboration, and independence of sources. A stronger score means that multiple pieces of evidence point in the same direction and that the claim is supported by a stable evidentiary pattern rather than by a single isolated source alone.

\begin{equation}
c = \mathrm{clip}\Big(\beta_0 + \beta_1\,\mathrm{ovl}_{\max} + \beta_2\,n_{\mathrm{extra}}
+ \beta_3\,n_{\mathrm{indep}} - \gamma\,\mathbb{1}[\text{self-only}],\; c_{\min},\; c_{\max}\Big)
\end{equation}

The exact coefficients are implementation-dependent, but the thesis-level interpretation is clear: confidence should reflect the strength and breadth of evidence. This is a valuable design choice because it gives the user a compact yet meaningful signal about trust, especially when several candidate claims look plausible at first glance.

Traceability is equally important. Every verdict should retain the supporting snippets and the rationale used to reach it. That way, a reader can inspect the evidence behind a conclusion and understand why the system made a particular decision. This is especially useful in a thesis context because it makes the system easier to explain in a viva and easier to compare with alternative approaches.

\section{Scope and Non-Goals}

AI-Scientist is designed as a domain-general research assistant rather than a narrow subject-specific tool. The framework is meant to accept questions from multiple academic fields, and the verification pipeline remains consistent across domains. This is a strong design advantage because it shows that the method is not tied to one particular benchmark topic \cite{petroni2020kilt,thakur2021beir,gu2021pubmedbert}.

At the same time, the system is not intended to behave like a general-purpose chatbot. A scope guard prevents non-research queries from entering the verification pipeline. This keeps the system focused on the class of problems it is designed to solve and preserves the clarity of the thesis contribution.

The thesis also excludes several tempting but unnecessary claims:
\begin{itemize}[leftmargin=*]
    \item autonomous scientific discovery is not claimed,
    \item exhaustive literature coverage is not claimed,
    \item full-text deep reading of every paper is not required,
    \item and replacing human judgement is not the objective.
\end{itemize}

These non-goals strengthen the work because they keep the system well-scoped in a positive way: the framework is ambitious enough to be meaningful, but focused enough to be credible.

\section{Verification Requirements}

The problem definition imposes several operational requirements on the system:
\begin{itemize}[leftmargin=*]
    \item a claim may be marked supported only when evidence clearly aligns with it,
    \item contradiction must be reported explicitly when evidence conflicts,
    \item weak claims should produce insufficient evidence rather than overconfident conclusions,
    \item confidence must account for evidence strength and corroboration,
    \item and every verdict must remain traceable to its source snippets.
\end{itemize}

These requirements are the foundation of the thesis evaluation. They also explain why a claim-level metric suite is needed later on. A system that performs well under these requirements is doing more than generating text; it is performing disciplined scientific verification.

\section{Adaptive Reasoning}

One of the most valuable features of the thesis problem definition is adaptive reasoning. If a claim appears weak on the first pass, the system should be able to request more evidence and re-check the claim before finalizing the report. This makes AI-Scientist feel more like a careful researcher and less like a fixed one-shot summarizer.

The adaptive loop is deliberately bounded. It is meant to improve weak or uncertain claims, not to search indefinitely. That constraint keeps the system efficient while still allowing it to handle difficult claims in a principled way.

\begin{figure}[H]
\centering
\begin{tikzpicture}[
    node distance=0.95cm,
    scale=0.94,
    transform shape,
    box/.style={rectangle, rounded corners, draw=black, fill=blue!12, align=center, minimum width=4.55cm, minimum height=0.98cm},
    arrow/.style={-{Latex[length=3mm]}, thick}
]
\node[box] (a) {Initial claim and evidence check};
\node[box, below=of a] (b) {Confidence below threshold};
\node[box, below=of b] (c) {Targeted retrieval for weak claim verification};
\node[box, below=of c] (d) {Re-verification and updated verdict};
\node[box, below=of d] (e) {Traceable report to the user};
\draw[arrow] (a) -- (b);
\draw[arrow] (b) -- (c);
\draw[arrow] (c) -- (d);
\draw[arrow] (d) -- (e);
\end{tikzpicture}
\caption{Adaptive reasoning loop used to refine uncertain claims}
\label{fig:adaptive_reasoning_loop}
\end{figure}

\section{Chapter Summary}

This chapter formalized the AI-Scientist problem as a verification-first research task. It defined the task input and output, the verdict space, the layered evidence model, confidence and traceability requirements, scope boundaries, and the adaptive reasoning loop. These definitions give the later architecture and evaluation chapters a clear foundation. The next chapter uses these requirements to describe the system architecture in detail.
